% META: {"id":"TCOMPL-CORR-ME-004","chapitre":"TCOMPL-CORRELATION-CAUSALITE","type_objet":"methode","capacites_codes":["C4"],"methodes":["M4"],"mode_creation":"ex_nihilo","sources_inspiration":[],"status":"approved","origin":"needs_review","status_history":["needs_review","audited","accepted","approved"],"release_acceptance":"RELEASE_OWNER_BATCH_ACCEPTANCE","acceptance_closure_digest":"sha256:9b3ccf9a81c5520fbb7e03b7b2d3e7bf2057a02834b6d908bf3be5f49f3b553f","acceptance_receipt_digest":"sha256:4fad86f0282e0fa4e0419cca1ad6379ae7af5a280c04a4a90880f90a1c044242"}
% BEGIN-VERIFY
% from sympy import Rational
% f = lambda t: Rational(7, 5) * t + Rational(1, 2)
% assert f(Rational(5, 2)) == 4
% assert f(6) == Rational(89, 10)
% END-VERIFY
\begin{fichemethode}{M4}{Interpoler ou extrapoler avec un ajustement affine, avec regard critique}

\quandUtiliser{%
  Utiliser cette méthode lorsque l'énoncé demande d'ESTIMER une valeur a partir
  de la droite d'ajustement $y = ax + b$ : a l'interieur de la plage observee
  (interpolation) ou au-dela (extrapolation).
}

\pasApas{%
  \begin{enumerate}
    \item Écrire l'équation de l'ajustement $y = ax + b$ (fiche M2) et noter la
      plage des $x$ observes $[x_{\min}\,;\,x_{\max}]$.
    \item Substituer la valeur demandee $x_0$ : $y_0 = a x_0 + b$.
    \item Qualifier l'estimation : INTERPOLATION si
      $x_0 \in [x_{\min}\,;\,x_{\max}]$, EXTRAPOLATION sinon.
    \item Porter un regard critique sur une extrapolation : plus $x_0$ s'eloigne
      des donnees, moins l'estimation est fiable (le modèle peut cesser d'être
      valable : saturation, changement de regime...).
    \item Conclure avec une phrase qui précise l'unité, l'arrondi et la reserve
      de validite.
  \end{enumerate}
}

\exempleRedige{%
  Avec l'ajustement $y = 1{,}4x + 0{,}5$ obtenu sur des donnees ou
  $x \in [1\,;\,4]$, estimer $y$ pour $x_0 = 2{,}5$ puis pour $x_0 = 6$.

  \medskip
  \commentaireMarge{$2{,}5 \in [1\,;\,4]$ : interpolation.}%
  Pour $x_0 = 2{,}5$ : $y_0 = 1{,}4 \times 2{,}5 + 0{,}5 = 4$ — interpolation,
  estimation fiable si l'ajustement est bon.

  \medskip
  \commentaireMarge{$6 \notin [1\,;\,4]$ : extrapolation.}%
  Pour $x_0 = 6$ : $y_0 = 1{,}4 \times 6 + 0{,}5 = 8{,}9$ — extrapolation au-dela
  des donnees : l'estimation suppose que la tendance linéaire se prolonge, ce
  qui doit être justifie par le contexte.
}

\pieges{%
  \begin{itemize}
    \item \textbf{Piège 1.} Presenter une extrapolation lointaine comme une
      certitude : le modèle n'est valide que pres des donnees observees.
    \item \textbf{Piège 2.} Utiliser la droite pour predire dans le sens inverse
      ($x$ en fonction de $y$) sans refaire de regression.
    \item \textbf{Piège 3.} Oublier les unités ou sur-arrondir les coefficients
      avant la substitution (arrondir seulement le résultat final).
  \end{itemize}
}

\verifier{%
  Contrôler que l'estimation est coherente avec le nuage (ordre de grandeur), et
  toujours indiquer si l'on interpole ou si l'on extrapole.
}

\sEntrainer{\refExos{M4}}

\end{fichemethode}
